% Replaces the paper's previous "external validation unavailable" claim.
% Include at most one or two result figures after trimming another subsection
% if the ICDE main-text limit would otherwise be exceeded.
\subsection{Independent Validation Across Reference Systems}
\label{sec:externalvalidation}
The preceding experiments evaluate a pair of taxonomies assigned to a small human microbiome cohort. To assess whether the query-stability construction transfers beyond that cohort and classifier pair, we conducted a second experiment on a separately published \emph{gut-to-soil} amplicon dataset. The publicly released QIIME~2 sequence tabulation contains \textbf{41,328 identical amplicon sequence variant (ASV) identifiers} classified by both SILVA~144 and GTDB~r232, together with one shared pooled frequency vector comprising \textbf{21,227,600 reads}. Thus, a reference change does not alter the observation, the frequency or the denominator of a query. This experiment evaluates pooled taxon-abundance rankings rather than subject-level or phenotype-specific contrasts. The original study and dual-classifier artifact are available via the QIIME~2 community tutorial (9 October 2026) and dataset DOI 10.5281/zenodo.15390940.

\paragraph{Annotation agreement is resolution-dependent.}
Among the same pooled reads, identical nonempty literal labels under both classifiers account for 81.86\% at the phylum level, 44.89\% at the family level and only 26.71\% at the genus level. The remaining mass comprises disagreements and missing annotations (Fig.~\ref{fig:externalagreement}); matching literal labels does not by itself establish cross-ontology conceptual identity. The identity of the most abundant named phylum is invariant in the two observed complete reference assignments (\emph{Pseudomonadota}), whereas the observed top-three and top-five sets differ. The dominant family and the unfiltered dominant genus also differ. This is direct evidence that observational stability of a query depends on its taxonomic resolution.

\paragraph{A stable answer can be certified beyond the observed references.}
Applying the exact two-anchor budgeted margin operator to all 115 competing phylum labels yields a strictly positive minimum for \emph{Pseudomonadota} at every budget, including unrestricted independent binary choices for each of the 41,328 ASVs. The smallest margin is \textbf{2,891,234 reads} (13.62\% of the common pooled denominator), providing a model-conditional certificate of top-one dominance rather than merely an observed two-reference agreement. The alternative partition into 2,928 SILVA-family blocks also retains a positive bound. These are robustness results under declared assignment worlds; no biological probability distribution over switching events is inferred.

\paragraph{A constructive vulnerability at finer granularity.}
As a nomenclature sensitivity analysis, we exclude noninformative genus labels, including \texttt{Incertae\_Sedis}, \texttt{uncultured} and \texttt{--other}, consistently in both classifiers. Both observed references then rank \emph{Rhodococcus} first. Nevertheless, a counterfactual based at the complete GTDB assignment switches just \textbf{two} ASVs to their SILVA genus labels, assigning an additional 77,422 reads to \emph{Rhodanobacter}. Its pairwise margin relative to \emph{Rhodococcus} changes from +75,129 to \textbf{--2,293 reads}. The first one-ASV edit cannot reverse that pairwise ordering, while the two-edit construction does (Fig.~\ref{fig:externalradius}). No reversal occurs for any whole-SILVA-family-block assignment in this example; family coherence is an explicit sensitivity assumption, \emph{not} an externally validated law of reannotation. All paired-ASV frequency denominators remain fixed throughout.

\paragraph{Correctness and scope.}
The optimized budgeted margin was compared against exhaustive enumeration on 1,320 random small-world cases, with maximum absolute numerical discrepancy \(8.88\times10^{-16}\). The external dataset thereby validates input pairing, exact operator behavior and the distinction between observed agreement and certificate strength in a different microbiome setting. It does \emph{not} independently validate population-level biological findings, site-group contrasts, phylogenetic admissibility assumptions, or performance relative to a production incremental-view-maintenance engine. The exact archived source hash, parsing code, witness identifiers and derived data are supplied in the reproducibility supplement.

% Insert figures as needed and adjust labels to match the main document:
% \begin{figure}[t]\centering\includegraphics[width=\columnwidth]{Fig_external_1_taxonomic_agreement.pdf}
% \caption{Read-weighted literal taxonomic agreement on 41,328 shared ASVs; unassigned rank labels are shown separately.}\label{fig:externalagreement}\end{figure}
% \begin{figure}[t]\centering\includegraphics[width=\columnwidth]{Fig_external_3_genus_witness.pdf}
% \caption{Exact worst pairwise genus margin under two-anchor budgets: a two-ASV witness reverses the pooled-rank contrast only under the independent-switch model.}\label{fig:externalradius}\end{figure}
